%% EXEMPLO FICTÍCIO — Atividade 4: ensaio experimental e validação
\begin{atividade}{Ensaios de bancada da revisão B e sessões de validação com o atleta-piloto}{
  tipo         = ensaio,
  modalidade   = EXT,
  alternativa  = PES,
  horas        = 70,
  inicio       = 2024-09,
  fim          = 2025-02,
  projeto      = {NEBULA, estimulação elétrica funcional para ciclismo},
  papel        = {Responsável pela parte técnica},
  supervisor   = {Rafael Fictício, supervisor do NEBULA; Dra. Carla Fictícia, fisioterapeuta do Depto. Clínico},
  registros    = {NRO-PRO-005-901; NRO-PRO-003-901; NRO-PRO-003-903; NBL-135},
  competencias = {I, II, IV, VI, VII},
  disciplinas  = {Laboratório de Eletrônica; Instrumentação Eletrônica; Probabilidade e Estatística},
}

\begin{contexto}
Uma placa que funciona na bancada ainda não é um sistema que serve a quem vai usá-lo. O plano de
validação do NEBULA (\texttt{NRO-PRO-005-901}) pedia duas etapas depois da montagem: a verificação
da revisão B contra os requisitos elétricos da USRS, na bancada; e doze sessões de treino com o
atleta-piloto --- uma pessoa com lesão medular, voluntária, de fora da UFMG ---, para confirmar que o
sistema faz o que ela precisa: pedalar por mais tempo. As sessões são a parte da ação de extensão da
equipe em que a tecnologia volta para quem motivou o projeto.
\end{contexto}

\begin{contribuicao}
Escrevi as seções de bancada do plano de validação, conduzi os ensaios de bancada e respondi pela
parte técnica das doze sessões: antes de cada uma, conferi a impedância dos eletrodos e carreguei os
parâmetros prescritos; durante, registrei os dados e monitorei a segurança elétrica. Quem conduziu as
sessões, definiu os parâmetros de estimulação e esteve com o atleta foi a fisioterapeuta do
Departamento Clínico; eu não posicionei eletrodos nem ajustei nada sem a indicação dela.
\end{contribuicao}

\begin{desenvolvimento}
\fichatecnica{
  protocolo    = {Plano de validação \texttt{NRO-PRO-005-901}, seções 4 (bancada) e 6 (sessões)},
  instrumentos = {Osciloscópio de 100\,MHz e sonda de corrente de 50\,MHz (calibrados em 2024); câmera térmica; cargas resistivas de 0,1\,\%; multímetro de 6\,½ dígitos},
  local        = {LABBIO, Escola de Engenharia da UFMG},
  amostra      = {Duas placas; um atleta-piloto voluntário, em 12 sessões},
  etica        = {Protocolo aprovado pelo Comitê de Ética em Pesquisa da UFMG, CAAE 00000000.0.0000.0000 (exemplo); termo de consentimento assinado},
  relatorio    = {\texttt{NRO-PRO-003-901} (bancada) e \texttt{NRO-PRO-003-903} (sessões)},
}

\textbf{A bancada.} Cada canal das duas placas foi ensaiado em três cargas de 0,1\,\% (\SI{500}{\ohm},
\SI{1}{\kilo\ohm} e \SI{1,5}{\kilo\ohm}), em seis amplitudes, com cinco repetições
(Figura~\ref{fig:ex-bancada}). A corrente foi medida com a sonda, e a carga de cada fase, pela integral
do pulso no próprio osciloscópio; o desbalanço de carga é a diferença entre as fases dividida pela
carga de uma delas. A incerteza combinada da amplitude, somando a da sonda (1\,\%) e a do canal do
osciloscópio (1,5\,\%) em quadratura, é de 1,8\,\%.

\begin{figure}[htbp]
  \centering
  \resizebox{0.92\linewidth}{!}{% A montagem dos ensaios de bancada (exemplo fictício).
\begin{tikzpicture}[
    bloco/.style={draw=nro-tinta, line width=0.5pt, fill=white, rounded corners=1.2pt,
      minimum height=12mm, text width=31mm, align=center, font=\footnotesize, inner sep=2pt},
    seta/.style={-{Stealth[length=1.8mm]}, line width=0.55pt, draw=nro-tinta},
    x=1mm, y=1mm]
  \node[bloco, fill=nro-synapse!18] (est) at (0,20) {Estimulador\\{\scriptsize placa rev.\,B no invólucro}};
  \node[bloco] (carga) at (52,20) {Carga de teste\\{\scriptsize 500\,$\Omega$, 1\,k$\Omega$ e 1,5\,k$\Omega$ (0,1\,\%)}};
  \node[bloco] (osc) at (104,20) {Osciloscópio\\{\scriptsize 100\,MHz, 4 canais}};
  \node[bloco] (note) at (0,0) {Notebook\\{\scriptsize registro dos parâmetros}};
  \node[bloco] (cam) at (52,0) {Câmera térmica\\{\scriptsize estágio de saída}};
  \node[bloco] (mult) at (104,0) {Multímetro de 6\,½ dígitos\\{\scriptsize calibração de $R_S$}};
  \draw[seta] (est) -- node[above, font=\scriptsize] {canal 1 a 4} (carga);
  \draw[seta] (carga) -- node[above, font=\scriptsize, align=center] {sonda de\\corrente} (osc);
  \draw[seta] (note) -- node[left, font=\scriptsize] {USB} (est);
  \draw[seta, densely dashed] (cam) -- (est.south east);
  \draw[seta, densely dashed] (mult) -- (carga.south east);
\end{tikzpicture}
}
  \caption{A montagem dos ensaios de bancada. As linhas tracejadas são as medições auxiliares.}
  \fonte{Elaborado pela autora.}
  \label{fig:ex-bancada}
\end{figure}

\textbf{As sessões.} Doze sessões de cerca de uma hora, em oito semanas, no LABBIO. Antes de cada uma,
eu media a impedância dos eletrodos e, acima de \SI{1}{\kilo\ohm} --- o limite que saiu da
Atividade~2 ---, a fisioterapeuta trocava o eletrodo. O indicador de desempenho definido no plano foi o
tempo de pedalada contínua até a cadência cair 20\,\% abaixo da inicial. As quatro primeiras sessões
usaram o perfil de estímulo da revisão A; as oito seguintes, o perfil ajustado pela fisioterapeuta com
a faixa maior de largura de pulso que a revisão B permite. Os dados foram registrados sem nenhuma
informação que identifique o atleta.
\end{desenvolvimento}

\begin{resultados}
Na bancada, as duas placas atenderam a todos os requisitos (Tabela~\ref{tab:ex-bancada}). Na carga de
\SI{1,5}{\kilo\ohm}, a corrente ficou limitada a \SI{78}{\milli\ampere}, como a tensão de conformidade
prevê (\SI{118}{\volt}/\SI{1,5}{\kilo\ohm}): nessa carga, as amplitudes acima disso ficam fora da
faixa de operação, e o firmware da revisão C passou a bloqueá-las.

\begin{table}[htbp]
  \centering\small
  \caption{Verificação da revisão B na bancada (duas placas, quatro canais, cinco repetições).}
  \label{tab:ex-bancada}
  \begin{tabular}{|l|c|c|}\hline
    \rowcolor{nro-fundo}\rotulo{Requisito} & \rotulo{Critério de aceite} & \rotulo{Pior resultado} \\ \hline
    Exatidão da amplitude (10 a 100\,mA, até 1\,k$\Omega$) & $\le$\,2\,\% & 0,9\,\% \\ \hline
    Desbalanço de carga entre as fases & $\le$\,2\,\% & 0,8\,\% \\ \hline
    Largura do pulso (50 a 500\,µs) & $\le$\,2\,µs de erro & 0,6\,µs \\ \hline
    Frequência (20 a 50\,Hz) & $\le$\,0,5\,\% & 0,1\,\% \\ \hline
    Disparo do limite de corrente & 105 a 115\,mA & 108 a 111\,mA \\ \hline
  \end{tabular}
  \fonte{Relatório de execução \texttt{NRO-PRO-003-901}.}
\end{table}

Nas sessões, o tempo até a queda de 20\,\% da cadência passou de 11 minutos, em média, nas quatro
primeiras sessões, para 19 minutos nas quatro últimas. Não houve evento adverso. O limite de corrente
por hardware disparou duas vezes, ambas por eletrodo descolado durante a pedalada: a saída foi
desligada antes de qualquer desconforto, como o projeto prevê. O atleta pediu que a parada de
emergência ficasse ao alcance da mão direita --- mudança feita no invólucro.
\end{resultados}

\begin{evidencias}
\evidencia{Arquivo controlado}{NRO-PRO-005-901}{O plano de validação, com as seções de bancada que escrevi}{SOMA › Arquivos}
\evidencia{Arquivo controlado}{NRO-PRO-003-901}{O relatório dos ensaios de bancada da revisão B}{SOMA › Arquivos}
\evidencia{Arquivo controlado}{NRO-PRO-003-903}{O relatório das doze sessões, sem dados pessoais}{SOMA › Arquivos (classe confidencial)}
\evidencia{Atividade}{NBL-135}{O cartão da validação, atribuído a mim e concluído}{SOMA › Atividades › NEBULA}
\evidencia{Presença}{Check-ins de set./2024 a fev./2025}{A presença no LABBIO nos dias de ensaio e de sessão}{SOMA › Equipe › Presença (a pedido)}
\end{evidencias}

\begin{aprendizados}
A bancada pôs em prática o que Instrumentação Eletrônica ensina sobre incerteza, e o que Probabilidade e
Estatística ensina sobre repetição; o Laboratório de Eletrônica me deu o hábito do procedimento
escrito. O que só as sessões ensinaram: trabalhar ao lado de uma profissional de saúde, respeitando o
que é dela; que um protocolo aprovado por comitê de ética é compromisso, não burocracia; e que o
requisito mais importante das doze sessões veio do próprio atleta --- a posição do botão de parada.
\end{aprendizados}
\end{atividade}
